% Ported verbatim from paper/top3_joint_moment_bound.tex lines 54-175.
% Only cross-reference wording was adapted; no mathematics was altered.
For a finite simple undirected graph $G$ of order $n\ge3$, put
\[
 \Lambda_3(G)=\lambda_1(G)+\lambda_2(G)+\lambda_3(G),
\]
where the adjacency eigenvalues are ordered algebraically. The signs in
this sum are retained. The extremal problem asks how large $\Lambda_3(G)$
can be when only $n$ is prescribed. The spectral-sum problem appears in Ebrahimi's thesis~\cite{Ebrahimi},
whose Theorem~3.4.1 gives the explicit asymptotic coefficient
$\limsup_n\max_{|G|=n}\Lambda_3(G)/n\le3/\sqrt5<1.3417$;
general bounds for sums of leading eigenvalues were developed by
Mohar~\cite{Mohar}. More recently,
Sivashankar, Tang and Wakhare~\cite[Proposition 6.1]{STW} proved
\begin{equation}\label{stw-bound}
 \Lambda_3(G)\le\frac{3+\sqrt5}{4}n-3.
\end{equation}
Their argument uses the rank-three absolute projection constant
$\gamma(3)=(1+\sqrt5)/2$; a proof of the requisite projection bound is
provided by Der\k{e}gowska and Lewandowska~\cite{DL}.

The two quantities in that reduction are naturally coupled. If
$X\in\R^3$ satisfies $\E XX^T=I_3$ and $X'$ is an independent copy, they are
\[
 A=\E|X\cdot X'|,\qquad B=|\E X|^2.
\]
The separate estimates $A\le(1+\sqrt5)/2$ and $B\le1$ give
\eqref{stw-bound}. A better graph bound therefore requires a quantitative
restriction on their simultaneous size. The mixed projection constant
expressing this question was already introduced in~\cite[Section 6.2]{STW}.
We work directly with its isotropic-law formulation.

Define
\begin{equation}\label{main-constant}
 c=\frac{106106902139519}{107310156800000}
       +\frac{41\sqrt3}{240}-\frac{53}{9000000}
   =1.284673255723918505\ldots.
\end{equation}
\begin{theorem}\label{main}
Every finite simple unweighted graph $G$ on $n\ge3$ vertices satisfies
\begin{equation}\label{main-bound}
 \Lambda_3(G)\le cn-3.
\end{equation}
Moreover, every Borel probability law on $\R^3$ satisfying
$\E XX^T=I_3$ obeys
\begin{equation}\label{law-bound}
 \E|X\cdot X'|+|\E X|^2\le2c.
\end{equation}
Only a finite second moment is assumed; zero vectors and unbounded radii
are allowed.
\end{theorem}
The coefficient decrease from~\eqref{stw-bound} is
$0.024343738651028919\ldots$, with the constant term $-3$ unchanged.
This comparison concerns the stated universal order-only bound. It does
not compare all estimates involving additional graph parameters, nor
assert that $c$ is optimal. For completeness we record the full chain of
explicit coefficients: $3/\sqrt5=1.341640786\ldots$
from~\cite[Theorem~3.4.1]{Ebrahimi}, then $(3+\sqrt5)/4$
from~\cite[Proposition~6.1]{STW}, and finally $c$.

Two adjacent results frame the standing of~\eqref{main-bound}. For the
two-term sum $S_2=\lambda_1+\lambda_2$, Huang and Wei~\cite{HW} determine
the exact maximum together with its unique extremal graph, whose three part
sizes are governed by $n$ modulo $7$; Paulin~\cite{Paulin} determines the
closures of the rescaled pairs
$\bigl((\lambda_1+1)/n,(\lambda_2+1)/n\bigr)$ and
$\bigl((\lambda_2+1)/n,(\lambda_3+1)/n\bigr)$, giving
$\lambda_1+\lambda_2\le8n/7-2$, $\lambda_3\le n/3-1$ and
$\lambda_2+\lambda_3\le2n/3-2$. Neither settles the three-term sum, because
the reduction behind~\eqref{stw-bound} couples two quantities, and it is
their joint control that forces the different geometry analysed here.

One prior announcement bears directly on our lower bound, and we record it
here rather than in a footnote. The self-archived Zenodo
deposits~\cite{Newton1,Newton2} of Newton, dated 4~and 24~August 2026 and
therefore earlier than this work, introduce the same extremal constant
$\cstar$ under the name of a pole--pentagon law, together with the same
optimising weight, and assert the matching upper bound
$\Lambda_3(G)\le\cstar n-3$; the assertion is conditional by its own terms,
since the final positivity layer is a machine certificate whose
admissibility the author explicitly leaves to independent audit
(Remark~\ref{rem-newton}). We claim no priority for the constant.
Theorem~\ref{main} is unconditional and self-contained, and is weaker in
the coefficient by $c-\cstar=0.003105058851750628\ldots$.

We use graphons to pass to an attained maximum. For a symmetric measurable
kernel $W:[0,1]^2\to[0,1]$, let
\begin{equation}\label{sig-definition}
 \Sigma_k(W)=\sup_{\E FF^T=I_k}
       \iint W(x,y)F(x)\cdot F(y)\,dx\,dy,
 \qquad \mathcal M=\sup_W\Sigma_3(W).
\end{equation}
The supremum is over real orthonormal $k$-frames. Since $W\in L^2([0,1]^2)$,
the operator $T_W$ is Hilbert--Schmidt and therefore compact; it is also
self-adjoint, so its spectrum is real and accumulates only at zero, and
$\Sigma_k$ is the sum of the first $k$ positive eigenvalues, padded by zero
if there are fewer than $k$ of them. This convention applies only to
graphons, not to $\Lambda_3(G)$.
The maximum $\mathcal M$ is attained. For a maximizing graphon, choose an
orthonormal top-three eigenframe $F=(f_1,f_2,f_3)^T$, with $f_1\ge0$,
and call $F_\#dx$ its canonical law.

\begin{theorem}\label{structure-main}
Let $W$ attain $\mathcal M$, let $\sum_i p_i\delta_{v_i}$ ($p_i>0$) be its
canonical law, and let $L$ be the diagonal matrix of the three leading
eigenvalues of $W$. Then:
\begin{enumerate}[leftmargin=2.2em,itemsep=2pt]
\item[(a)] the support has at most $24$ points;
\item[(b)] $W$ agrees almost everywhere with the binary step kernel
$\mathbf1_{\{v_i\cdot v_j>0\}}$ on the corresponding positive-measure
fibers;
\item[(c)] the support points satisfy
\begin{gather*}
 \sum_i p_iv_iv_i^T=I_3,\qquad v_i^TLv_i=\mathcal M,\qquad v_{i1}>0,\\
 v_i\cdot v_j\ne0,\qquad
 B_{ij}=\mathbf1_{\{v_i\cdot v_j>0\}},\qquad
 \sum_jp_jB_{ij}v_j=Lv_i;
\end{gather*}
\item[(d)] for every positive-dot clique $J$ in the support, the matrices
$\{v_iv_i^T:i\in J\}$ are linearly independent; in particular $|J|\le6$.
\end{enumerate}
\end{theorem}
This assertion concerns every original maximizer, not merely the
existence of a maximizer with small support.
